\begin{frame}{역사: Cover--Hart(1967)의 보증}
  1967년 Cover와 Hart의 논문(``Nearest Neighbor Pattern
  Classification'')이 이론을 정립했다.
  \vskip0.8em
  \begin{block}{믿기 힘든 정리}
    ``가장 가까운 이웃의 답을 따른다''는 단순한 방법도
    \textbf{데이터가 무한히 많으면} 오류율이
    \kb{베이즈 오류율}(어떤 방법으로도 줄일 수 없는 이론적 하한)의
    \textbf{2배를 넘지 못한다}.
  \end{block}
  \begin{itemize}
    \item 더 놀라운 점: \textbf{1차원}(특징 1개)에서는 2배가 아니라
      \textbf{정확히 베이즈 오류율 자체}로 수렴
    \item 그래서 kNN은 ``맹목적인 방법''이 아니라,
      \textbf{어디까지 틀릴 수 있는지 이론적으로 보증된 방법}
  \end{itemize}
\end{frame}
